\section{Problema detectado y/o faltante}

En la actualidad, gran parte de los clubes de fútbol en Argentina administran a su masa societaria e instalaciones mediante sistemas de gestión tradicionales e infraestructuras tecnológicas fragmentadas. Esta situación genera ineficiencias operativas crónicas y pérdidas económicas significativas que, con frecuencia, pasan desapercibidas para las comisiones directivas. Entre las principales problemáticas se destacan la alta tasa de rotación de socios (abandono o morosidad), la evasión en los controles de acceso los días de partido (mediante el préstamo de carnets físicos o el uso de capturas de pantalla de códigos QR estáticos) y el desaprovechamiento comercial de los espacios alquilables de las instituciones.

Asimismo, se detecta una clara ausencia de herramientas integrales en el mercado actual orientadas al sector deportivo local que permitan:

\begin{itemize}
    \item Identificar y prevenir la morosidad societaria de manera proactiva mediante el análisis de datos, evitando la desafiliación antes de que ocurra.
    \item Garantizar un control de accesos robusto, seguro y ágil en escenarios de alta concurrencia y nula conectividad a internet.
    \item Automatizar la atención al socio y la gestión de cobros para reducir la cantidad desproporcionada de horas-hombre destinadas a tareas administrativas repetitivas.
    \item Optimizar, digitalizar y rentabilizar el uso de los espacios compartidos y servicios adicionales que el club ofrece a su comunidad.
\end{itemize}

\textbf{SocioUnido} nace para resolver este faltante, funcionando como un ecosistema unificado y proactivo que combina la gestión administrativa integral, la retención inteligente de socios mediante IA y la ciberseguridad en accesos físicos en una única plataforma.
